\documentclass{article}
\usepackage[margin=1in]{geometry}
\usepackage{amsmath, amssymb, amsthm, enumitem, multicol, booktabs}

\everymath{\displaystyle}
\newtheorem*{theorem*}{Theorem}
\newtheorem*{definition*}{Definition}
\newcommand{\vct}[1]{\langle #1 \rangle}
\newcommand{\vu}{\mathbf{u}}
\newcommand{\vv}{\mathbf{v}}
\newcommand{\vw}{\mathbf{w}}
\newcommand{\vx}{\mathbf{x}}
\newcommand{\vn}{\mathbf{n}}
\newcommand{\va}{\mathbf{a}}
\newcommand{\vb}{\mathbf{b}}

\newcommand{\colvecTwo}[2]{
  \begin{bmatrix}
    #1 \\
    #2
  \end{bmatrix}
}
\newcommand{\colvecThree}[3]{
  \begin{bmatrix}
    #1 \\
    #2 \\
    #3
  \end{bmatrix}
}

\newcommand{\casesTwo}[2]{
  \begin{cases}
    #1 \\
    #2
  \end{cases}
}
\newcommand{\casesThree}[3]{
  \begin{cases}
    #1 \\
    #2 \\
    #3
  \end{cases}
}

\title{Linear Algebra: A Modern Introduction}
\author{David Poole}
\date{}

\begin{document}
\maketitle
\setcounter{section}{2}
\section{Matrices}
\subsection{Matrix Operations}
\begin{definition*}
  A \textbf{\textit{matrix}} is a rectangular array of numbers called the \textbf{\textit{entries}},
  or \textbf{\textit{elements,}} of the matrix.
\end{definition*}
\subsubsection*{Exercises}
\begin{enumerate}
    \setcounter{enumi}{16}
  \item $
    \begin{bmatrix}
      0 & 0 \\
      1 & 0
    \end{bmatrix}$
    \setcounter{enumi}{18}
  \item $
    \begin{bmatrix}
      650.00 & 462.50 \\
      675.00 & 406.25
    \end{bmatrix}$
    \setcounter{enumi}{22}
  \item $
    \begin{bmatrix}
      2\va_1+\va_2-\va_3 & 3\va_1-\va_2+6\va_3 & \va_2+4\va_3
    \end{bmatrix}$
  \item $
    \begin{bmatrix}
      \vb_1-2\vb_2 \\
      -3\vb_1+\vb_2+\vb_3 \\
      2\vb_1-\vb_3
    \end{bmatrix}$
  \item $
    \begin{bmatrix}
      2 & 3 & 0 \\
      -6 & -9 & 0 \\
      4 & 6 & 0
    \end{bmatrix}
    +
    \begin{bmatrix}
      0 & 0 & 0 \\
      1 & -1 & 1 \\
      0 & 0 & 0
    \end{bmatrix}
    +
    \begin{bmatrix}
      2 & -12 & -8 \\
      -1 & 6 & 4 \\
      1 & -6 & -4
    \end{bmatrix}
    $
    \setcounter{enumi}{30}
  \item $
    \begin{bmatrix}
      3 & 2  & 0 \\
      -1 & 1 & 0 \\
      0 & 0 & 5
    \end{bmatrix}
    $
\end{enumerate}

\subsection{Matrix Algebra}
\begin{theorem*}
  $(AB)^T=B^TA^T$
\end{theorem*}
\begin{proof}
  If $A$ is $m\times n$ and $B$ is $n\times r$, then $B^T$ is $r\times n$ and $A^T$ is $n\times m$.
  Thus, the product $B^TA^T$ is defined and is $r\times m$. Since $AB$ is $m\times r$, $(AB)^T$ is
  $r\times m$, and so $(AB)^T$ and $B^TA^T$ have the same size. We must now prove that their
  corresponding entries are equal.

  We denote the $i$th row of a matrix $X$ by $\operatorname{row}_i(X)$ and its $j$th column by $
  \operatorname{col}_j(X)$. Using these conventions, we see that
  \begin{align*}
    [(AB)^T]_{ij} & = (AB)_{ji} \\
    & = \operatorname{row}_j(A)\cdot\operatorname{col}_i(B) \\
    & = \operatorname{col}_i(B)\cdot\operatorname{row}_j(A) \\
    & = \operatorname{row}_i(B)^T\cdot\operatorname{col}_j(A)^T \\
    & = (B^TA^T)_{ij}.
  \end{align*}
\end{proof}
\subsubsection*{Exercises}
\begin{enumerate}
  \item $
    \begin{bmatrix}
      5 & 4 \\
      3 & 5
    \end{bmatrix}
    $
    \setcounter{enumi}{4}
  \item $B=2A_1+A_2$
    \setcounter{enumi}{8}
  \item $
    \begin{bmatrix}
      x & y \\
      -5x+2y & -x+y
    \end{bmatrix}
    $
    \setcounter{enumi}{12}
  \item Yes.
    \setcounter{enumi}{22}
  \item $
    \begin{bmatrix}
      a & b \\
      0 & a
    \end{bmatrix}
    $
    \setcounter{enumi}{45}
  \item $\sum_{i}\sum_{j}a_{ij}^2$
\end{enumerate}

\subsection{The Inverse of a Matrix}
\end{document}
